% ============================================================
% Relatório técnico -- cp-02 Arquiteturas (Tópicos Avançados em IA)
%
% COMO USAR NO OVERLEAF (template Springer LNCS):
%   1. Abra o projeto que veio do template LNCS.
%   2. Substitua TODO o conteúdo do samplepaper.tex por este arquivo.
%      (O llncs.cls e o splncs04.bst do template continuam lá; não mexa.)
%   3. Suba as 5 imagens (aba de arquivos, "Upload"), só os arquivos,
%      sem as pastas:
%        dados/amostras.png
%        comparacao/curvas_treino.png
%        comparacao/matriz_confusao.png
%        comparacao/curva_roc.png
%        comparacao/discordancias.png
%   4. Compile.
%
% Todo número deste texto sai de um arquivo versionado no repositório;
% a Tabela "Rastreabilidade" (seção 4) diz qual.
% ============================================================
\documentclass[runningheads]{llncs}

\usepackage[T1]{fontenc}
\usepackage[utf8]{inputenc}
\usepackage[brazil]{babel}
\usepackage{graphicx}
\usepackage{booktabs}
\usepackage{hyperref}
\renewcommand\UrlFont{\color{blue}\rmfamily}

\begin{document}

\title{Menos é mais: enxugando a ResNet-18 do MedMNIST v2 para a
escala do PneumoniaMNIST}
\titlerunning{Enxugando a ResNet-18 do MedMNIST v2}

\author{Davi Pires Aquino de Carvalho}
\authorrunning{D. P. A. de Carvalho}
\institute{Universidade Federal Rural de Pernambuco (UFRPE)\\
Tópicos Avançados em Inteligência Artificial}

\maketitle

\begin{abstract}
Reproduzimos o \textit{baseline} ResNet-18 do benchmark MedMNIST v2
(2023) no PneumoniaMNIST, a partir do código oficial dos autores, e
propomos uma modificação de arquitetura que remove o último estágio
residual, reduz os blocos por estágio de dois para um e divide a
largura por dois. A rede cai de 11,2 milhões para 307 mil parâmetros.
Treinadas com a mesma receita, a rede original obteve AUC 0,949 e
acurácia 0,872 no teste (publicado: 0,944 e 0,854), e a rede
modificada obteve AUC 0,964 e acurácia 0,875, treinando 6,8 vezes mais
rápido. As curvas de treino mostram que a rede original memoriza o
conjunto de treino e tem a perda de validação crescendo 78\% após o
seu mínimo, contra 15\% na rede modificada.

\keywords{CNN \and ResNet \and MedMNIST \and Pneumonia \and Arquitetura}
\end{abstract}

\section{Introdução}

Redes convolucionais projetadas para bases grandes, como o ImageNet, são
frequentemente reaproveitadas sem ajuste em bases muito menores. A
ResNet-18~\cite{he2016} tem cerca de 11 milhões de parâmetros; o
PneumoniaMNIST~\cite{medmnistv2}, uma das bases do benchmark MedMNIST v2,
tem 4.708 imagens de treino de $28\times28$ pixels e apenas duas classes.
Nessa desproporção, é razoável suspeitar que boa parte da capacidade da
rede não contribui para generalizar e só serve para memorizar o treino.

O objetivo desta prática, dentro da disciplina, é demonstrar como uma
modificação de arquitetura impacta uma rede neural. Para isso
(i) selecionamos um artigo de 2023 que aplica uma CNN a uma base pública,
(ii) reproduzimos o seu resultado com o código oficial dos autores,
(iii) modificamos a arquitetura de forma substancial,
(iv) retreinamos a rede modificada na mesma base, com a mesma receita, e
(v) comparamos as duas quantitativa e qualitativamente.

A hipótese testada é que a ResNet-18 do artigo está superdimensionada
para o PneumoniaMNIST, e que uma versão com uma fração dos parâmetros
alcança desempenho equivalente.

\section{Artigo e rede base}

O artigo escolhido é o MedMNIST v2~\cite{medmnistv2}, publicado na
\textit{Scientific Data} em 2023. Ele padroniza 12 bases de imagens
biomédicas 2D em $28\times28$ pixels e publica \textit{baselines} com
código aberto (\texttt{github.com/MedMNIST/experiments}). A rede base é a
ResNet-18 desse código, que é a variante para imagens pequenas adaptada do
repositório \texttt{kuangliu/pytorch-cifar}: a primeira convolução é
$3\times3$ com passo 1 e não há \textit{max-pooling} inicial, de modo que
a imagem entra nos blocos residuais em resolução cheia.

Para o PneumoniaMNIST, o artigo reporta AUC 0,944 e acurácia 0,854 para a
ResNet-18 em resolução 28, e AUC 0,948 e acurácia 0,854 para a
ResNet-50, que tem cerca de duas vezes mais parâmetros. A ResNet-50
não melhorar a acurácia foi o primeiro indício que motivou a modificação.

\section{Método}

\subsection{Dados}

O PneumoniaMNIST deriva de 5.856 radiografias de tórax pediátricas de
Kermany et~al.~\cite{kermany2018}, recortadas e reduzidas para
$1\times28\times28$ pelos autores do MedMNIST, com rótulo binário
(normal ou pneumonia). Usamos a divisão oficial: 4.708 imagens de treino,
524 de validação e 624 de teste. O arquivo foi conferido pelo MD5
registrado no código dos autores antes de qualquer treino.

A Fig.~\ref{fig:amostras} mostra exemplos das duas classes e a contagem
por divisão. A base é desbalanceada: 74,2\% das imagens de treino e de
validação são de pneumonia. O teste tem proporção diferente, 62,5\%
(390 de 624): os autores do MedMNIST tiraram treino e validação do
conjunto de treino da base de origem, e usaram como teste o conjunto de
validação dela.

\begin{figure}[htbp]
    \centering
    \includegraphics[width=\textwidth]{amostras.png}
    \caption{Exemplos de radiografias de treino de cada classe (esquerda) e
    número de imagens por classe em cada divisão (direita).}
    \label{fig:amostras}
\end{figure}

\subsection{Protocolo}

Copiamos o projeto oficial dos autores (commit \texttt{70b6b3a}) e
alteramos apenas a arquitetura. O diff toca dois arquivos: em
\texttt{models.py}, a classe \texttt{ResNet} passou a aceitar a largura e
o número de estágios como parâmetros; em
\texttt{train\_\allowbreak and\_\allowbreak eval\_\allowbreak pytorch.py}, uma linha permite selecionar a nova
rede. Verificamos que, com os valores padrão, a ResNet-18 continua
idêntica à original: mesmos nomes e formatos de parâmetros, mesma contagem
e mesma saída para os mesmos pesos. Otimizador, \textit{scheduler},
pré-processamento, divisão dos dados e métricas são os dos autores, de
modo que a arquitetura é a única variável entre as duas execuções.

A receita de treino é a do artigo: Adam~\cite{kingma2015}, taxa de
aprendizado 0,001, lote de 128, 100 épocas, com a taxa dividida por 10
nas épocas 50 e 75. O script dos autores guarda o
\textit{checkpoint} da época com maior AUC de validação, e é esse que
avaliamos no teste. As duas redes foram treinadas em CPU de 4 núcleos, uma
após a outra, para que o tempo de uma não fosse inflado pela disputa de
processador com a outra.

\subsection{A modificação}

A Tabela~\ref{tab:params} mostra onde estão os parâmetros da ResNet-18
nesta configuração. O quarto estágio concentra 75,2\% da rede e opera
sobre mapas de $4\times4$, resolução em que resta pouca estrutura
espacial de uma radiografia de $28\times28$.

\begin{table}[htbp]
    \centering
    \caption{Distribuição dos parâmetros da ResNet-18 do artigo
    (1 canal de entrada, 2 classes).}
    \label{tab:params}
    \begin{tabular}{lrrc}
        \toprule
        Parte & Parâmetros & \% da rede & Mapa de saída \\
        \midrule
        conv. inicial + estágio 1 & 148.672 & 1,3 & $64\times28\times28$ \\
        estágio 2 & 525.568 & 4,7 & $128\times14\times14$ \\
        estágio 3 & 2.099.712 & 18,8 & $256\times7\times7$ \\
        \textbf{estágio 4} & \textbf{8.393.728} & \textbf{75,2} & $\mathbf{512\times4\times4}$ \\
        classificador & 1.026 & 0,0 & -- \\
        \midrule
        total & 11.168.706 & 100 & \\
        \bottomrule
    \end{tabular}
\end{table}

A rede modificada, chamada aqui de \textit{enxuta}, aplica três cortes
(Tabela~\ref{tab:arq}): remove o quarto estágio, usa um bloco residual por
estágio em vez de dois e divide por dois o número de filtros de cada
estágio. O bloco residual e a estrutura geral da ResNet são mantidos.

\begin{table}[htbp]
    \centering
    \caption{Arquitetura antes e depois da modificação.}
    \label{tab:arq}
    \begin{tabular}{lcc}
        \toprule
        & ResNet-18 (artigo) & Enxuta \\
        \midrule
        Estágios residuais & 4 & 3 \\
        Blocos por estágio & 2, 2, 2, 2 & 1, 1, 1 \\
        Filtros por estágio & 64, 128, 256, 512 & 32, 64, 128 \\
        Resolução ao longo da rede & $28\to28\to14\to7\to4$ & $28\to28\to14\to7$ \\
        Parâmetros & 11.168.706 & 307.042 \\
        \bottomrule
    \end{tabular}
\end{table}

A justificativa tem três partes: a proporção de cerca de 2.400
parâmetros por imagem de treino; a própria tabela do artigo, em que
dobrar a rede não melhorou a acurácia; e o fato de o estágio que mais pesa
ser justamente o que enxerga a menor resolução.

\subsection{Etapas, ferramentas e uso de IA}

O trabalho seguiu as etapas: (1) busca e seleção do artigo, restrita a
publicações de 2023 em diante; (2) cópia do código oficial e verificação
de que a rede base não foi alterada; (3) análise da distribuição de
parâmetros e escolha da modificação; (4) treino das duas redes no
PneumoniaMNIST; (5) comparação e redação deste relatório. As bibliotecas
são as do código dos autores (PyTorch, \texttt{medmnist}) mais
\texttt{matplotlib}, \texttt{scikit-learn} e \texttt{tensorboard} para as
figuras.

Todo o código, a execução dos experimentos e o rascunho deste relatório
foram produzidos pelo \textbf{Claude Code} (Anthropic), sob a direção do
autor, conforme a orientação da disciplina de conduzir o trabalho
comandando a IA. O autor definiu as restrições (artigo de 2023 em diante,
base de tema próximo), escolheu a base e a modificação entre as opções
propostas pela IA e revisou os resultados. Para evitar alucinações, cada
número do texto foi conferido contra o arquivo que o produziu
(Tabela~\ref{tab:rastro}), e as referências foram conferidas nas fontes
dos próprios autores.

\section{Resultados}

\subsection{Reprodução da rede base}

A ResNet-18 obteve AUC 0,949 e acurácia 0,872 no teste, contra 0,944 e
0,854 publicados. A diferença é compatível com a variação entre execuções:
o script dos autores não fixa semente aleatória. Consideramos a reprodução
válida e usamos essa execução como referência do ``antes''.

\subsection{Comparação quantitativa}

\begin{table}[htbp]
    \centering
    \caption{Resultados no conjunto de teste (624 imagens). Sensibilidade é
    a fração de casos de pneumonia detectados; especificidade, a fração de
    casos normais reconhecidos como normais.}
    \label{tab:resultados}
    \begin{tabular}{lccc}
        \toprule
        & ResNet-18 (artigo) & Enxuta & Publicado~\cite{medmnistv2} \\
        \midrule
        Parâmetros & 11.168.706 & \textbf{307.042} & -- \\
        AUC & 0,949 & \textbf{0,964} & 0,944 \\
        Acurácia & 0,872 & \textbf{0,875} & 0,854 \\
        Sensibilidade & 0,985 & 0,982 & -- \\
        Especificidade & 0,684 & 0,697 & -- \\
        Treino (100 épocas) & 119 min & \textbf{17 min} & -- \\
        Inferência por imagem & 4,3--4,8 ms & \textbf{0,55--0,61 ms} & -- \\
        \bottomrule
    \end{tabular}
\end{table}

A Tabela~\ref{tab:resultados} resume o resultado. Com 36 vezes menos
parâmetros, a rede enxuta obteve AUC 0,015 maior e acurácia praticamente
igual, treinou 6,8 vezes mais rápido e infere de 7 a 8 vezes mais rápido
(três medições). A hipótese se confirma: a redução não custou desempenho.

\subsection{Curvas de treino}

A Fig.~\ref{fig:curvas} mostra a perda e a AUC de validação por época.
As duas redes levam a perda de treino a valores entre $10^{-6}$ e
$10^{-5}$ e chegam a 100\% de acurácia no treino, ou seja, ambas
memorizam o conjunto de treino, o que é esperado de redes profundas com
capacidade suficiente~\cite{zhang2017}. A diferença está na validação: a
perda de validação da ResNet-18 atinge o mínimo na época 6 (0,070) e
termina 78\% acima dele (0,125), enquanto a da rede enxuta termina apenas
15\% acima do seu mínimo (de 0,076 para 0,088). É o padrão de
sobreajuste que motivou a modificação, e ele é bem mais fraco na rede
menor.

Os \textit{checkpoints} escolhidos pela AUC de validação vieram das
épocas 25 (ResNet-18) e 10 (enxuta), antes das reduções da taxa de
aprendizado. Nas duas redes, as últimas 50 épocas não melhoraram a
validação.

\begin{figure}[htbp]
    \centering
    \includegraphics[width=\textwidth]{curvas_treino.png}
    \caption{Perda (escala logarítmica) e AUC de validação por época. Linhas
    contínuas: treino; tracejadas: validação. Os círculos marcam o
    \textit{checkpoint} usado no teste.}
    \label{fig:curvas}
\end{figure}

\subsection{Comparação qualitativa}

As matrizes de confusão (Fig.~\ref{fig:confusao}) mostram que as duas
redes erram do mesmo lado: detectam 98\% dos casos de pneumonia, mas
reconhecem só cerca de 70\% dos casos normais (160 e 163 de 234),
classificando os demais como pneumonia. É o reflexo do desbalanceamento
do treino, e a modificação não o altera.

\begin{figure}[htbp]
    \centering
    \includegraphics[width=\textwidth]{matriz_confusao.png}
    \caption{Matrizes de confusão no teste. Cor: fração da linha; rótulo:
    contagem e fração.}
    \label{fig:confusao}
\end{figure}

As curvas ROC (Fig.~\ref{fig:roc}) ficam próximas, com a da rede enxuta
acima na região de poucos falsos positivos, que é onde a diferença de AUC
se concentra. Imagem a imagem, as duas redes erram juntas em 60 radiografias
e discordam em apenas 38: a enxuta acerta 20 delas e a ResNet-18, 18
(Fig.~\ref{fig:discordancias}). A diferença de acurácia entre as duas é,
portanto, de duas imagens.

\begin{figure}[htbp]
    \centering
    \includegraphics[width=0.55\textwidth]{curva_roc.png}
    \caption{Curvas ROC das duas redes no teste.}
    \label{fig:roc}
\end{figure}

\begin{figure}[htbp]
    \centering
    \includegraphics[width=\textwidth]{discordancias.png}
    \caption{Exemplos de radiografias em que as redes discordam, com o
    rótulo verdadeiro e a previsão e confiança de cada uma.}
    \label{fig:discordancias}
\end{figure}

\subsection{Rastreabilidade}

A Tabela~\ref{tab:rastro} indica de qual arquivo do repositório sai cada
número apresentado nesta seção, para que qualquer um deles possa ser
conferido ou regenerado.

\begin{table}[htbp]
    \centering
    \caption{Origem de cada número, no repositório do trabalho
    (pasta \texttt{code/cnn-pneumonia-medmnist/}).}
    \label{tab:rastro}
    \small
    \begin{tabular}{ll}
        \toprule
        Número & Arquivo \\
        \midrule
        AUC e acurácia & \texttt{resultados/*/pneumoniamnist\_log.txt} \\
        Matriz de confusão, discordâncias & CSVs de teste em \texttt{resultados/} \\
        Curvas, épocas escolhidas & \texttt{resultados/*/Tensorboard\_Results/} \\
        Tempo de treino & \texttt{resultados/execucao.log} \\
        Tempo de inferência & \texttt{resultados/comparacao.txt} \\
        Parâmetros por estágio & \texttt{verificacao/contar\_parametros.py} \\
        Distribuição das classes & \texttt{dados/pneumoniamnist.npz} \\
        \bottomrule
    \end{tabular}
\end{table}

\section{Discussão e limitações}

\textbf{Execução única.} Cada rede foi treinada uma vez, sem semente fixa.
Como a diferença de acurácia é de duas imagens, não afirmamos que a rede
enxuta é melhor; afirmamos que ela não perdeu desempenho. Repetir cada
treino algumas vezes e comparar médias permitiria uma conclusão mais
forte sobre a AUC.

\textbf{Mudança de proporção entre validação e teste.} A validação tem
74,2\% de pneumonia e o teste, 62,5\%. Como as redes tendem a responder
``pneumonia'' na dúvida, erram mais onde há mais casos normais, o que é
compatível com a queda de acurácia de 96--97\% na validação para
cerca de 87\% no teste nas duas redes.

\textbf{Especificidade baixa.} As duas redes reconhecem só cerca de 70\%
dos casos normais. Isso decorre dos dados, não da arquitetura, e
exigiria outras intervenções, como ponderar as classes na função de perda.

\textbf{Escala das imagens.} As radiografias originais chegam a
$2.916\times2.713$ pixels; em $28\times28$, nada aqui indica desempenho
clínico. O objeto de estudo é a arquitetura, não o diagnóstico.

\textbf{Espaço para reduzir mais.} A rede enxuta também chega a 100\% de
acurácia no treino, o que sugere que ela ainda tem capacidade sobrando.

\section{Conclusão}

Reproduzimos o \textit{baseline} ResNet-18 do MedMNIST v2 no
PneumoniaMNIST, com resultado próximo ao publicado, e mostramos que a rede
está superdimensionada para essa base: o seu maior estágio, com 75\% dos
parâmetros, pode ser removido junto com metade dos blocos e metade dos
filtros. A rede resultante tem 36 vezes menos parâmetros, treina 6,8 vezes
mais rápido, sofre menos sobreajuste e mantém o desempenho no teste
(AUC 0,964 contra 0,949; acurácia 0,875 contra 0,872). Para esta base, o
tamanho da rede do artigo não se traduz em ganho de generalização.

\begin{thebibliography}{6}

\bibitem{he2016}
He, K., Zhang, X., Ren, S., Sun, J.: Deep residual learning for image
recognition. In: Proceedings of the IEEE Conference on Computer Vision
and Pattern Recognition (CVPR), pp. 770--778 (2016)

\bibitem{kermany2018}
Kermany, D.S., Goldbaum, M., et al.: Identifying medical diagnoses and
treatable diseases by image-based deep learning. Cell \textbf{172}(5),
1122--1131.e9 (2018)

\bibitem{kingma2015}
Kingma, D.P., Ba, J.: Adam: A method for stochastic optimization. In:
International Conference on Learning Representations (ICLR) (2015)

\bibitem{medmnistv2}
Yang, J., Shi, R., Wei, D., Liu, Z., Zhao, L., Ke, B., Pfister, H.,
Ni, B.: MedMNIST v2 -- A large-scale lightweight benchmark for 2D and 3D
biomedical image classification. Scientific Data \textbf{10}(1), 41
(2023)

\bibitem{zhang2017}
Zhang, C., Bengio, S., Hardt, M., Recht, B., Vinyals, O.: Understanding
deep learning requires rethinking generalization. In: International
Conference on Learning Representations (ICLR) (2017)

\end{thebibliography}

\end{document}
